\newpage

\section[FUNDAMENTAÇÃO TEÓRICA]{Fundamentação Teórica}

2. FUNDAMENTAÇÃO TEÓRICA
2.1 SOLO - nesse subitem você pode definir o solo (produto do intemperismo sobre as rochas) e fazer comentários sobre sua origem (solos residuais, solos transportados e as características dos horizontes), fazer comentários sobre o tamanho das partículas, comentários sobre o comportamento desse material com a variação da umidade (limites de consistência)
2.2 SOLO COMO MATERIAL CONSTRUTIVO - nesse subitem você pode retomar os benefícios dessas técnicas construtivas, tanto em aspectos ambientais (disponibilidade do material, reciclabilidade, consumo de energia, geração de resíduos) como em relação ao conforto dos usuários (térmico, acústico, etc.). Aqui você pode citar brevemente algumas técnicas como taipa de pilão, cob, pau a pique, bloco de terra comprimida e o próprio adobe. 
2.2.1. Adobe - nesse subitem você pode se aprofundar um pouco mais na técnica do adobe: origem, regiões do planeta em que sua utilização é mais comum, tipo de solo que melhor se adapta à técnica, adições comuns para o melhoramento de suas propriedades (pelos e gordura animal, adição de fibras vegetais, etc) e estratégias para a durabilidade das construções (beirais, impermeabilização das fundações), etc. Nesse subitem foque em práticas tradicionais. O uso de produtos industrializados como cal e cimento serão abordados no próximo subitem.
2.3 TÉCNICAS DE MELHORAMENTO DE SOLOS - nesse subitem você pode contextualizar a necessidade de melhoramento de solos para que ele usado para fins construtivos, você pode trazer os principais tipos de estabilização: mecânica, granulométrica/física e química.
2.3.1 Solo-cimento - nesse subitem você pode trazer mais detalhes sobre o solo-cimento, principais características, tipos de cimento comumente utilizados para esse fim, principais reações químicas, tipos de solo que melhor se adaptam à técnica e principais limitações (tendência de trincamento, perda de resistência quando exposto à umidade excessiva, etc.).

% A Fundamentação Teórica ou Revisão da Literatura deve contextualizar teoricamente o problema e apresentar o estágio atual de conhecimento sobre o assunto. A revisão de literatura é a parte que sustenta todo o projeto, pois é por meio dela que se encontram indicações para a solução do problema identificado. Devido a sua importância, ela deve ser atual, abrangente e com profundidade. Por outro lado, deve deter-se a assuntos específicos do projeto de pesquisa em desenvolvimento, sem incluir aspectos desnecessários e de forma exaustiva. Trata-se da apresentação do embasamento teórico sobre o qual se fundamentará o trabalho. Durante a escrita, o autor deve lembrar de que o texto não deve ser construído apenas com uma sequência de recortes e citações de outros autores sobre o tema. Segundo Gil (2018), a revisão bibliográfica não é constituída apenas por referências ou sínteses do relato de estudos, mas por discussão crítica das obras citadas. 

% ORIENTAÇÕES GERAIS Escreva todo seu texto utilizando estilo normal, fonte Times New Roman, tamanho 12, exceto as citações com mais de três linhas, notas de rodapé, paginação, legendas e fontes das figuras e tabelas, que devem ser em tamanho menor e uniforme. Todo os elementos prétextuais e textuais deverão ser formatados para um tamanho de página A4 (210 x 297 mm), limitado por margens esquerda e superior de 3 cm e direita e inferior de 2 cm. Todo o texto deve ser digitado com espaçamento 1,5 entre as linhas e alinhamento justificado, exceto as citações de mais de três linhas, notas de rodapé, referências, legendas das figuras e tabelas. As referências bibliográficas, ao final do trabalho, devem ser separadas entre si por um espaço simples em branco. Os títulos das seções primárias devem começar em página ímpar (anverso) e ser separados do texto que os sucede por um espaço entre as linhas de 1,5. Do mesmo modo, os títulos das subseções devem ser separados do texto que os procede e que os sucede por um espaço entre as linhas de 1,5. Lista de figuras, lista de tabelas, lista de abreviaturas e siglas, lista de símbolos, sumário, referências, apêndice(s) e anexo(s) são títulos sem indicativo numérico e devem ser centralizados. 287 As folhas ou páginas pré-textuais devem ser contadas, mas não numeradas. Os trabalhos devem ser digitados ou datilografados somente no anverso. Todas as folhas, a partir da folha de rosto, devem ser contadas sequencialmente, considerando somente o anverso (a capa não é contada - página “zero”). A numeração deve aparecer, a partir da primeira folha da parte textual (Introdução), em algarismos arábicos, no canto superior direito da folha. 

% Tabelas As tabelas devem ser enumeradas sequencialmente, citadas no texto, inseridas o mais próximo possível do trecho a que se referem e padronizadas conforme o Instituto Brasileiro de Geografia e Estatística (IBGE, 1993). • Toda tabela que ultrapassar a dimensão da página em número de linhas e tiver poucas colunas, pode ter o centro apresentado em duas ou mais partes, lado a lado, na mesma página, separando-se as partes por um traço vertical duplo e repetindo-se o cabeçalho; • Toda tabela que ultrapassar a dimensão da página em número de colunas, e tiver poucas linhas, pode ter o centro apresentado em duas ou mais partes, uma abaixo da outra, na mesma página, repetindo-se o cabeçalho das colunas indicadoras e os indicadores de linha; • Para toda tabela que ultrapassar as dimensões da página: cada página deve ter o conteúdo do topo e o cabeçalho da tabela ou o cabeçalho da parte; cada página deve ter uma das seguintes indicações: continua para a primeira, conclusão para a última e continuação para as demais. O título das tabelas deve ser incluído na linha imediatamente anterior à tabela e centralizado, utilizando fonte Times New Roman, tamanho 10 e cor preta, conforme Tabela 1. Para configurar o título das tabelas acesse o menu “Referências” e clique em “Inserir Legenda”. A fonte deve ser identificada logo abaixo, mesmo que de autoria própria, indicando “acervo particular” ou “autoria própria”.


% Equações e fórmulas Buscando facilitar a leitura, as equações e fórmulas devem ser destacadas no texto e, se necessário, numeradas com algarismos arábicos entre parênteses, alinhados à direita. Todas as variáveis envolvidas nas equações bem como a unidade do parâmetro calculado devem ser explicitadas ao longo do texto ou em seguida à apresentação da equação. A Equação 1 apresenta a fórmula para o cálculo da velocidade escalar média, como forma de exemplificação. vm = ∆S ∆T (1) Onde: vm: velocidade escalar média (m/s); ΔS: distância total (m); ΔT: intervalo de tempo (s). 

% Ilustrações Qualquer que seja o tipo de ilustração, sua identificação aparece na parte superior, precedida da palavra designativa (desenho, fluxograma, gráfico, mapa, planta, figura, imagem, entre outros), seguida de seu número de ordem de ocorrência no texto, em algarismos arábicos, e do respectivo título. A fonte deve ser identificada logo abaixo, mesmo que de autoria própria, indicando “acervo particular” ou “autoria própria”. A ilustração deve ser citada no texto e inserida o mais próximo possível do trecho a que se refere. A Figura 1 mostra uma representação esquemática de um solo com estrutura alveolar, com a intenção de demonstrar o processo de configuração de figuras.

\begin{figure}[H]
    \centering
    \caption{Teste}
    \label{fig:placeholder}
    \includegraphics[width=0.5\linewidth]{image.jpg}
    
    {\small Fonte: Elaborado pelo autor (2026). \par}
\end{figure}

% Citações As referências a autores ou transcrição de informações retiradas de outras fontes devem seguir as diretrizes da ABNT NBR 10520. As citações diretas consistem na transcrição textual de parte da obra do autor consultado, e as citações indiretas consistem na elaboração de um texto baseado na obra do autor consultado. Nas citações, as chamadas pelo sobrenome do autor devem ser em letras maiúsculas e minúsculas e, quando estiverem entre parênteses, devem ser em letras maiúsculas. Exemplo 1: De acordo com Miguez, Veról e Rezende (2015), o processo de urbanização gera grandes modificações no ambiente natural, alterando os padrões de uso e ocupação do solo e agravando os problemas de enchentes. Exemplo 2: A mistura entre agregado reciclado cimentício e agregado reciclado de cerâmica vermelha pode favorecer a formação de compostos cimentantes, uma vez que a fração fina do resíduo de cerâmica vermelha tem desempenho de material pozolânico (SILVA, 2014). As citações diretas, no texto, de até três linhas, devem estar contidas entre aspas duplas. As citações diretas, no texto, com mais de três linhas, devem ser destacadas com recuo de 4 cm da margem esquerda, com letra menor que a do texto utilizado e sem as aspas. Nas citações diretas deve-se especificar no texto: a(s) página(s), volume(s), tomo(s) ou seção(ões) da fonte consultada. Estas informações devem seguir o ano da publicação, separado(s) por vírgula e precedido(s) pelo termo que o(s) caracteriza, de forma abreviada. Nas citações indiretas, a indicação da(s) página(s) consultada(s) não é obrigatória. Exemplo de citação direta de até três linhas: “Cimento, no sentido geral da palavra pode ser descrito como um material com propriedades adesivas e coesivas que o fazem capaz de unir fragmentos minerais na forma de uma unidade compacta.” (NEVILLE, 2015, p.1). Exemplo de citação direta com mais de três linhas: Uma geogrelha é definida como um material polimérico (isto é, geossintético) composto por conjuntos paralelos de arestas elásticas conectadas com aberturas do tamanho suficiente para permitir a penetração do solo, pedras ou outros materiais geotécnicos do entorno. As geogrelhas normalmente são feitas com polietileno de alta densidade (HDPE) e polipropileno (PP). A principal função de uma geogrelha é o reforço. (DAS; SOBHAN, 2019, p. 662).